% e-Shop 系統需求規格書（Phase 1）
% 編譯：tectonic srs.tex
\documentclass[12pt,a4paper]{article}

\usepackage[a4paper,top=2.5cm,bottom=2.5cm,left=2.5cm,right=2.5cm]{geometry}
\usepackage{fontspec}
\usepackage{xeCJK}
\setmainfont{Times New Roman}
\setCJKmainfont{DFKai-SB}[AutoFakeBold=2.5]
\xeCJKsetup{CJKecglue={ }}
\linespread{1.3}

\usepackage{titlesec}
\usepackage{tocloft}
\usepackage{xltabular}
\usepackage{needspace}
\usepackage{array}
\usepackage[table]{xcolor}
\usepackage{enumitem}
\usepackage{tikz}
\usetikzlibrary{positioning,arrows.meta,fit,backgrounds,calc}
\usepackage{amsmath}
\usepackage{hyperref}
\hypersetup{hypertexnames=false,colorlinks=true,linkcolor=black,urlcolor=blue,citecolor=black}
\usepackage{xurl}

% ---------- 標題格式（依 SRS.pdf 範本） ----------
\titleformat{\section}{\centering\large\bfseries}{Section \thesection}{1em}{}
\titleformat{\subsection}{\bfseries}{\thesubsection}{0.5em}{}
\titleformat{\subsubsection}{\normalfont}{\thesubsubsection}{0.5em}{}
\titlespacing*{\section}{0pt}{0pt}{1.5em}
\titlespacing*{\subsection}{0pt}{1.2em}{0.5em}
\titlespacing*{\subsubsection}{0pt}{0.8em}{0.3em}

\renewcommand{\contentsname}{\centering 目錄 (Table of Contents)}
\renewcommand{\cftsecpresnum}{Section }
\setlength{\cftsecnumwidth}{5.5em}
\renewcommand{\cftsecleader}{\cftdotfill{\cftdotsep}}
\renewcommand{\cftsecfont}{\bfseries}
\renewcommand{\cftsecpagefont}{\bfseries}
\setlength{\cftsubsecindent}{1.5em}
\setlength{\cftsubsubsecindent}{3em}
\setlength{\cftsubsecnumwidth}{2.2em}
\setlength{\cftsubsubsecnumwidth}{3em}

\renewcommand{\figurename}{圖}
\renewcommand{\tablename}{表}

% ---------- 需求表格 ----------
\definecolor{headgray}{gray}{0.88}
\setlist{nosep,leftmargin=2em}

% 需求編號｜需求描述｜優先順序
\newenvironment{reqtable}{%
  \xltabular{\linewidth}{|>{\centering\arraybackslash}p{2.6cm}|X|>{\centering\arraybackslash}p{1.9cm}|}
  \hline\rowcolor{headgray}\textbf{需求編號} & \textbf{需求描述} & \textbf{優先順序}\\\hline\endhead
}{\endxltabular}
\newcommand{\req}[3]{#1 & #2 & #3\\\hline}

% 編號｜描述（兩欄）
\newenvironment{twocol}[2]{%
  \xltabular{\linewidth}{|>{\centering\arraybackslash}p{2.6cm}|X|}
  \hline\rowcolor{headgray}\textbf{#1} & \textbf{#2}\\\hline\endhead
}{\endxltabular}
\newcommand{\row}[2]{#1 & #2\\\hline}

% 待確認標記
\newcommand{\assume}{\textsuperscript{\textbf{[假設]}}}

% ---------- 組員資料：請自行填入 ----------
\newcommand{\projectname}{FUWA FUWA e-Shop 線上寵物用品商城}
\newcommand{\projectnameEN}{FUWA FUWA e-Shop}
\newcommand{\docdate}{10/07/2026}

\begin{document}

% ================= 封面 =================
\begin{titlepage}
\centering
\vspace*{1.5cm}
{\Large\bfseries \projectname\\[0.3em] (\projectnameEN)\par}
\vspace{1.5cm}
{\LARGE\bfseries 系統需求規格書\par}
\vspace{0.5cm}
{\Large\bfseries Software Requirements Specification (SRS)\par}
\vspace{0.5cm}
{\large\bfseries Version: 1.0\par}
\vspace{1.2cm}
\renewcommand{\arraystretch}{1.4}
\begin{tabular}{ccc}
\textbf{姓名} & \textbf{學號} & \textbf{E-mail}\\
徐法恩 & 113370239 & t113370239@ntut.org.tw\\
魏提安 & 113590032 & t113590032@ntut.org.tw\\
鄧玉璿 & 113590029 & t113590029@ntut.org.tw\\
黃柏瑋 & 113590023 & t113590023@ntut.org.tw\\
陳俊佑 & 113AC1012 & t113AC1012@ntut.org.tw\\
\end{tabular}
\vfill
{\large\bfseries Department of Computer Science \& Information Engineering\\
National Taipei University of Technology\par}
\vspace{1cm}
{\large\bfseries \docdate\par}
\vspace*{1cm}
\end{titlepage}

% ================= 目錄 =================
\pagenumbering{roman}
\setcounter{page}{1}
\tableofcontents
\clearpage
\pagenumbering{arabic}
\setcounter{page}{1}

% ================= Section 1 =================
\section{簡介 (Introduction)}

\subsection{目的 (Purpose)}
本文件為資料庫系統課程專題 Phase 1 的系統需求規格書，描述「\projectname」（以下簡稱 e-Shop）的系統需求。e-Shop 是一個線上電子商務系統，讓顧客可以在網路上搜尋、選購寵物用品並查詢自己的訂單，讓工作人員可以管理商品、處理訂單、訂定折扣政策並產生銷售報表，讓平台管理員可以管理使用者帳號、角色與權限。

本專題的重點在資料庫設計，因此本文件除了列出系統功能外，也特別整理系統需要儲存的資料（2.4 節）與商業規則及完整性限制（2.8 節），作為 Phase 2 繪製 ER Diagram 與設計資料庫 schema 的依據。

本系統主要目標為：
\begin{itemize}
  \item 顧客可以搜尋商品、使用購物車選購商品、結帳並查詢歷史訂單。
  \item 系統於結帳時自動計算小計、稅、運費，並套用運費、季節性與特殊活動三種折扣。
  \item 工作人員可以管理商品與折扣政策、更新訂單狀態，並查看日、週、月銷售報表。
  \item 平台管理員可以管理使用者帳號、指派角色與權限，並停用或啟用帳號。
  \item 以關聯式資料庫的主鍵、外鍵、唯一性與檢查約束確保資料的正確性與一致性。
\end{itemize}

\subsection{系統名稱 (Identification)}
本專案範圍包含建置下列主系統與各子系統。主系統為：

\textbf{FUWA FUWA 線上商城系統 (FUWA FUWA e-Shop System, FES)}

各子系統分別為：
\begin{itemize}
  \item 帳號與權限子系統 (Account and Authorization Subsystem, AAS)
  \item 商品管理子系統 (Product Management Subsystem, PMS)
  \item 購物車子系統 (Shopping Cart Subsystem, SCS)
  \item 訂單管理子系統 (Order Management Subsystem, OMS)
  \item 折扣管理子系統 (Discount Management Subsystem, DMS)
  \item 報表統計子系統 (Report and Statistics Subsystem, RSS)
  \item 商店管理子系統 (Store Management Subsystem, STS)
  \item 資料庫子系統 (Database Subsystem, DBS)
\end{itemize}

\subsection{概觀 (Overview)}
本系統採用前後端分離的三層式架構。前端以 React 18、TypeScript 與 Vite 開發單頁式網頁應用程式（SPA），同一個網站依登入角色顯示三種介面：顧客購物介面、工作人員管理介面與平台管理介面。後端以 Python 的 FastAPI 框架提供 REST API，透過 SQLAlchemy 存取資料庫。資料庫選用 PostgreSQL，原因是它支援完整的 SQL 標準、交易與各種完整性約束。

本文件的組織如下：Section 1 說明文件目的與系統名稱；Section 2 描述系統架構、操作概念、功能性需求、資料需求、非功能性需求、介面需求、其他需求，以及商業規則與完整性限制；最後為名詞解釋、參考資料與附錄。

\subsection{符號描述 (Notation Description)}
各系統元件與需求以下列編號表示：

\begin{twocol}{編號}{說明}
\row{FES 1.0.0}{The FES system will be labeled with the number 1.0.0.}
\row{AAS 1.1.n}{The AAS components will be labeled with the number 1.1.n.}
\row{PMS 1.2.n}{The PMS components will be labeled with the number 1.2.n.}
\row{SCS 1.3.n}{The SCS components will be labeled with the number 1.3.n.}
\row{OMS 1.4.n}{The OMS components will be labeled with the number 1.4.n.}
\row{DMS 1.5.n}{The DMS components will be labeled with the number 1.5.n.}
\row{RSS 1.6.n}{The RSS components will be labeled with the number 1.6.n.}
\row{STS 1.7.n}{The STS components will be labeled with the number 1.7.n.}
\row{DBS 1.8.n}{The DBS components will be labeled with the number 1.8.n.}
\row{XXX-F-nnn}{子系統 XXX 的功能性需求 (Functional Requirements)}
\row{FES-N-nnn}{系統的非功能性需求 (Non-Functional Requirements)}
\row{US-nn}{使用者故事 (User Story)}
\row{DR-nn}{資料需求 (Data Requirement)}
\row{BR-nn}{商業規則與完整性限制 (Business Rule / Integrity Constraint)}
\end{twocol}

優先順序分為兩級：「必要」為課程題目以 shall 規定，或本組決定一定要完成的需求；「選配」為題目標示 optional 或 should、時間允許才實作的需求。

\clearpage
% ================= Section 2 =================
\section{系統(System)}

\subsection{系統描述 (System Description)}
本系統分為八個子系統：帳號與權限子系統（AAS 1.1.0）負責登入、登出、註冊與角色授權；商品管理子系統（PMS 1.2.0）負責商品搜尋與商品資料的新增、查詢、修改、刪除；購物車子系統（SCS 1.3.0）負責購物車的加入、修改與移除；訂單管理子系統（OMS 1.4.0）負責結帳、確認付款、建立訂單、查詢歷史訂單、更新訂單狀態，以及取消訂單與退貨申請；折扣管理子系統（DMS 1.5.0）負責折扣政策的管理與結帳時的折扣計算；報表統計子系統（RSS 1.6.0）負責日、週、月銷售報表；商店管理子系統（STS 1.7.0）負責商店資料與商店評價；資料庫子系統（DBS 1.8.0）負責所有資料的儲存與完整性限制。

系統有三種使用者角色：顧客（Customer）、工作人員（Staff）與平台管理員（Administrator），一位使用者可以同時具有多種角色，例如平台管理員兼任工作人員。初始資料庫中至少有一位平台管理員，工作人員與平台管理員角色由平台管理員指派。未登入的訪客只能瀏覽與搜尋商品，在介面上歸在顧客這一側。目前工作人員可以管理所有商店的資料；讓各商店的工作人員只管理自家商店的獨立權限屬於未來擴充，不在本期範圍內。

\subsubsection{系統架構圖 (System Architecture Diagram)}
圖 1 為本系統的架構。使用者以瀏覽器開啟前端 SPA；三種角色使用同一個網站，前端依登入角色顯示顧客購物介面、工作人員管理介面或平台管理介面，不是三套獨立系統。前端以 HTTP(S) 傳送 JSON 格式的 REST API 請求給後端；後端依請求呼叫對應的子系統，再透過 SQLAlchemy 以 SQL 存取 PostgreSQL 資料庫。寄送訂單確認信為選配功能，以虛線表示。

\begin{figure}[htbp]
\centering
\begin{tikzpicture}[
  font=\small,
  box/.style={draw,rounded corners=2pt,align=center,minimum height=0.9cm,fill=white},
  sub/.style={draw,align=center,minimum width=1.85cm,minimum height=0.95cm,fill=white,font=\footnotesize},
  actor/.style={align=center,font=\small},
  arr/.style={-{Stealth},thick},
  darr/.style={{Stealth}-{Stealth},thick}
]
% 使用者
\node[actor] (cus) at (2.9,0) {\textbf{顧客}＋訪客\\(Customer / Guest)};
\node[actor] (stf) at (7.6,0) {\textbf{工作人員}\\(Staff)};
\node[actor] (adm) at (12.3,0) {\textbf{平台管理員}\\(Administrator)};
% 前端
\draw[rounded corners=3pt,fill=yellow!10] (0,-1.1) rectangle (15.2,-3.3);
\node[box,minimum width=4.2cm] (fe1) at (2.9,-1.95) {顧客購物介面};
\node[box,minimum width=4.2cm] (fe2) at (7.6,-1.95) {工作人員管理介面};
\node[box,minimum width=4.2cm] (fe3) at (12.3,-1.95) {平台管理介面};
\node[font=\footnotesize] at (7.6,-2.95) {瀏覽器（Browser）：同一網站依登入角色顯示，React 18 + TypeScript + Vite（SPA）};
\draw[arr] (cus) -- (fe1);
\draw[arr] (stf) -- (fe2);
\draw[arr] (adm) -- (fe3);
% 後端
\draw[rounded corners=3pt,fill=blue!5] (0,-4.3) rectangle (15.2,-7.5);
\node[font=\small\bfseries] at (7.6,-4.65) {後端伺服器：FastAPI（REST API）};
\foreach \i/\c/\t in {0/AAS/帳號權限,1/PMS/商品,2/SCS/購物車,3/OMS/訂單,4/DMS/折扣,5/RSS/報表,6/STS/商店}
  \node[sub] at ({1.3+\i*2.1},-5.55) {\c\\\t};
\node[box,minimum width=7.2cm,minimum height=0.8cm] at (4.25,-6.85) {JWT 驗證、角色授權、bcrypt 密碼雜湊};
\node[box,minimum width=5.6cm,minimum height=0.8cm] at (11.6,-6.85) {SQLAlchemy（ORM / SQL）};
\draw[darr] (7.6,-3.3) -- node[right,font=\footnotesize]{HTTP(S)，REST / JSON} (7.6,-4.3);
% 資料庫與郵件
\node[box,minimum width=4.4cm,minimum height=1.2cm,fill=green!8] (db) at (7.6,-8.9) {\textbf{DBS}\\PostgreSQL 資料庫};
\draw[darr] (7.6,-7.5) -- node[right,font=\footnotesize]{SQL} (db.north);
\node[box,dashed,minimum width=3.6cm,minimum height=1.2cm] (mail) at (13.0,-8.9) {SMTP 郵件伺服器\\（選配）};
\draw[arr,dashed] (13.0,-7.5) -- node[right,font=\footnotesize]{確認信} (mail.north);
\end{tikzpicture}
\caption{e-Shop 系統架構圖}
\end{figure}

\subsection{操作概念 (Operational Concepts or User Stories)}

\textbf{Scenario 1：訪客操作概念 (Guest Operational Concepts)}

訪客經由顧客購物介面瀏覽商品，可以使用商品管理子系統（PMS）以關鍵字搜尋商品、查看商品詳細資料，也可以透過商店管理子系統（STS）查看商店資訊。訪客若要使用購物車或下單，必須先透過帳號與權限子系統（AAS）註冊為顧客並登入。

\textbf{Scenario 2：顧客操作概念 (Customer Operational Concepts)}

顧客透過 AAS 以 loginid 與密碼登入後，可以搜尋商品並將商品加入購物車子系統（SCS），在購物車中修改數量或移除商品。結帳時，訂單管理子系統（OMS）會請折扣管理子系統（DMS）計算目前有效的折扣，並顯示小計、折扣、稅、運費與總額。顧客確認付款後，系統建立訂單、扣除庫存並清空購物車，訂單狀態為 received。之後顧客可以在 OMS 查詢歷史訂單與每筆訂單目前的狀態。訂單仍為 received 時，顧客可以直接取消訂單；商店已開始處理（processing）、已出貨（shipping）或已完成（closed）後，顧客只能提出退貨申請，可以選擇要退的商品與數量，須經工作人員同意才會退貨，且每筆訂單只能申請一次。訂單完成後，顧客可以對商店評分與評論（選配）。

\textbf{Scenario 3：工作人員操作概念 (Staff Operational Concepts)}

工作人員透過 AAS 登入後進入工作人員管理介面，可以使用 PMS 新增、查詢、修改與刪除商品及庫存；使用 OMS 查詢顧客的訂單，將訂單狀態依序更新為 processing、shipping、closed，並審核顧客的退貨申請；使用 DMS 新增、查詢、修改與刪除運費、季節性與特殊活動三種折扣政策；使用報表統計子系統（RSS）查看各商店的日、週、月銷售報表。

\textbf{Scenario 4：平台管理員操作概念 (Administrator Operational Concepts)}

平台管理員透過 AAS 登入後進入平台管理介面，可以新增使用者帳號、查詢與修改使用者資料、替使用者指派或移除角色（管理角色與權限），以及停用或啟用帳號。若平台管理員同時具有工作人員身分，也可以使用 Scenario 3 的所有功能。

表 1 以使用者故事的形式整理各角色的主要操作。

\begin{xltabular}{\linewidth}{|>{\centering\arraybackslash}p{1.4cm}|>{\centering\arraybackslash}p{1.8cm}|X|}
\caption{使用者故事 (User Stories)}\\
\hline\rowcolor{headgray}\textbf{編號} & \textbf{角色} & \textbf{使用者故事}\\\hline\endfirsthead
\hline\rowcolor{headgray}\textbf{編號} & \textbf{角色} & \textbf{使用者故事}\\\hline\endhead
US-01 & 訪客 & 身為訪客，我想以關鍵字搜尋商品，以便快速找到想買的東西。\\\hline
US-02 & 訪客 & 身為訪客，我想自行註冊為顧客，以便使用購物車與下單。（選配）\\\hline
US-03 & 顧客 & 身為顧客，我想以 loginid 與密碼登入與登出，以便安全地使用我的帳號。\\\hline
US-04 & 顧客 & 身為顧客，我想將商品加入購物車並調整數量或移除，以便決定要購買的內容。\\\hline
US-05 & 顧客 & 身為顧客，我想在結帳時看到小計、折扣、稅、運費與總額，以便確認要付多少錢。\\\hline
US-06 & 顧客 & 身為顧客，我希望有效期內的折扣自動套用，以便不會錯過優惠。\\\hline
US-07 & 顧客 & 身為顧客，我想在確認付款後完成下單，並收到訂單確認信（選配）。\\\hline
US-08 & 顧客 & 身為顧客，我想查詢我的歷史訂單與訂單目前的狀態。\\\hline
US-09 & 顧客 & 身為顧客，我想在商店尚未處理訂單（received）時直接取消訂單。\\\hline
US-10 & 顧客 & 身為顧客，我想對商店已處理、已出貨或已完成的訂單選擇要退的商品與數量並提出退貨申請，並查詢處理結果。\\\hline
US-11 & 顧客 & 身為顧客，我想對購買過的商店評分與留下評論。（選配）\\\hline
US-12 & 工作人員 & 身為工作人員，我想新增、查詢、修改、刪除商品與庫存，以便維持正確的商品資訊。\\\hline
US-13 & 工作人員 & 身為工作人員，我想查看新進訂單並依序更新為 processing、shipping、closed。\\\hline
US-14 & 工作人員 & 身為工作人員，我想審核顧客的退貨申請，決定同意或拒絕。\\\hline
US-15 & 工作人員 & 身為工作人員，我想新增運費、季節性與特殊活動折扣，並由系統指派唯一折扣代碼。\\\hline
US-16 & 工作人員 & 身為工作人員，當我設定的特殊活動折扣期間與同商品既有折扣重疊時，我希望系統拒絕並提示我。\\\hline
US-17 & 工作人員 & 身為工作人員，我想查看各商店的日、週、月銷售報表（日期、商店名稱、訂單數、銷售總額）。\\\hline
US-18 & 工作人員 & 身為工作人員，我想查看各商店的暢銷商品與指定期間的購買統計。（選配）\\\hline
US-19 & 平台管理員 & 身為平台管理員，我想建立與管理使用者帳號、管理使用者的角色與權限，並能停用違規的帳號，以便維持平台正常運作。\\\hline
\end{xltabular}

\needspace{12\baselineskip}
\subsection{功能性需求 (Functional Requirements)}

\needspace{8\baselineskip}\textbf{帳號與權限子系統 (AAS)}
\begin{reqtable}
\req{AAS-F-001}{使用者可以使用 loginid 與密碼登入系統。}{必要}
\req{AAS-F-002}{使用者可以登出系統。}{必要}
\req{AAS-F-003}{系統依使用者角色（顧客、工作人員、平台管理員）授權可使用的功能；一位使用者可以同時具有多種角色，例如平台管理員兼任工作人員。}{必要}
\req{AAS-F-004}{訪客可以自行註冊為顧客，註冊時須填寫 loginid、姓名、密碼、Email 與地址，loginid 不可與既有使用者重複。}{選配}
\req{AAS-F-005}{平台管理員可以新增使用者帳號。}{必要}
\req{AAS-F-006}{平台管理員可以查詢與修改使用者的姓名、Email 與地址，並可以停用或啟用帳號；平台管理員不能停用自己的帳號。}{必要}
\req{AAS-F-007}{使用者可以查看與修改自己的個人資料及密碼。}{選配}
\req{AAS-F-008}{平台管理員可以替使用者指派或移除角色（顧客、工作人員、平台管理員），以管理各使用者可使用的功能；系統至少須保留一位啟用中的平台管理員。}{必要}
\end{reqtable}

\needspace{8\baselineskip}\textbf{商品管理子系統 (PMS)}
\begin{reqtable}
\req{PMS-F-001}{訪客與顧客可以用關鍵字搜尋商品名稱與描述。}{必要}
\req{PMS-F-002}{訪客與顧客可以依分類、商店篩選商品，並依價格或上架時間排序。}{選配}
\req{PMS-F-003}{系統顯示商品詳細資料，包含名稱、描述、圖片、價格、庫存、所屬商店，以及目前適用的特殊活動折扣。}{必要}
\req{PMS-F-004}{工作人員可以新增商品，須指定所屬商店、名稱、描述、價格、庫存與圖片。}{必要}
\req{PMS-F-005}{工作人員可以查詢商品，包含已下架的商品。}{必要}
\req{PMS-F-006}{工作人員可以修改商品資料與庫存數量。}{必要}
\req{PMS-F-007}{工作人員可以刪除商品；若商品已被訂單引用，系統改為將商品下架而不刪除。}{必要}
\req{PMS-F-008}{商品頁顯示「買了這個商品的人也買了」推薦。}{選配}
\end{reqtable}

\needspace{8\baselineskip}\textbf{購物車子系統 (SCS)}
\begin{reqtable}
\req{SCS-F-001}{顧客可以將商品加入購物車並指定數量。}{必要}
\req{SCS-F-002}{顧客可以修改購物車中商品的數量。}{必要}
\req{SCS-F-003}{顧客可以從購物車移除商品。}{必要}
\req{SCS-F-004}{購物車顯示每項商品的單價、數量、小計，以及預估的折扣、稅、運費與總額。}{必要}
\req{SCS-F-005}{購物車內容儲存在資料庫中，顧客登出再登入或換裝置後仍然保留。}{必要}
\req{SCS-F-006}{顧客輸入的數量不是正整數或超過目前庫存時，系統拒絕並提示。}{必要}
\end{reqtable}

\needspace{8\baselineskip}\textbf{訂單管理子系統 (OMS)}
\begin{reqtable}
\req{OMS-F-001}{顧客結帳時，系統在伺服器端計算商品小計、各項折扣、稅、運費與訂單總額，並顯示金額明細。}{必要}
\req{OMS-F-002}{顧客結帳時可以確認或修改收件人姓名、電話與地址，預設帶入顧客的個人資料。}{必要}
\req{OMS-F-003}{顧客確認付款後，系統在同一個資料庫交易內再次檢查庫存、扣除庫存、建立訂單與訂單明細、記錄套用的折扣並清空購物車；任一步驟失敗則全部回復，不建立訂單。}{必要}
\req{OMS-F-004}{新建立的訂單狀態為 received，並記錄下單時間。}{必要}
\req{OMS-F-005}{顧客可以查詢自己的歷史訂單列表與每筆訂單的明細（品名、單價、數量、小計、折扣、稅、運費、總額、狀態）。}{必要}
\req{OMS-F-006}{工作人員可以查詢所有訂單，並依狀態與日期篩選。}{必要}
\req{OMS-F-007}{工作人員可以將訂單狀態依序更新為 processing、shipping、closed，系統記錄每次狀態變更的時間與操作的工作人員；訂單有待審核的退貨申請時不可更新狀態。}{必要}
\req{OMS-F-008}{顧客可以查看訂單的狀態歷程（各狀態的變更時間）。}{選配}
\req{OMS-F-009}{顧客下單成功後，系統自動寄送訂單確認 Email。}{選配}
\req{OMS-F-010}{顧客可以取消狀態為 received 的訂單，系統立即將訂單狀態改為 cancelled 並回補庫存。}{必要}
\req{OMS-F-011}{顧客可以對狀態為 processing、shipping 或 closed 的訂單提出退貨申請，選擇要退的商品與數量並填寫退貨原因；每筆訂單只能申請一次。}{必要}
\req{OMS-F-012}{工作人員可以查詢退貨申請，並同意或拒絕。同意時，系統回補退貨商品的庫存、依 BR-26 計算並記錄退款金額；若訂單的所有商品都已退貨，訂單狀態改為 returned，否則維持原狀態。拒絕時須填寫拒絕原因，訂單維持原狀態。}{必要}
\req{OMS-F-013}{顧客可以查詢自己退貨申請的處理狀態與結果。}{必要}
\end{reqtable}

\needspace{8\baselineskip}\textbf{折扣管理子系統 (DMS)}
\begin{reqtable}
\req{DMS-F-001}{工作人員可以新增折扣政策，須指定折扣類型（運費、季節性、特殊活動）、描述、開始與結束日期，以及該類型所需的規則參數；系統自動指派唯一的折扣代碼。}{必要}
\req{DMS-F-002}{運費折扣的規則為：訂單金額大於或等於門檻金額時免運費。}{必要}
\req{DMS-F-003}{季節性折扣的規則為：在指定期間內，整筆訂單金額乘上折扣率作為折抵金額。}{必要}
\req{DMS-F-004}{特殊活動折扣的規則為折扣率或「買 X 送 Y」，並須指定適用的商品（一個或多個）。}{必要}
\req{DMS-F-005}{工作人員可以查詢折扣政策，並依類型與期間篩選。}{必要}
\req{DMS-F-006}{工作人員可以修改折扣政策；折扣代碼不可修改。}{必要}
\req{DMS-F-007}{工作人員可以刪除折扣政策；已被訂單使用過的政策不可刪除，只能提前結束或停用。}{必要}
\req{DMS-F-008}{設定特殊活動折扣的適用商品時，若與同一商品既有特殊活動折扣的期間重疊，系統拒絕並提示衝突的折扣代碼。}{必要}
\req{DMS-F-009}{結帳時系統自動套用有效期內所有符合條件的折扣，顧客不需輸入折扣代碼；運費、季節性與特殊活動三種類型可同時套用於同一筆訂單，同一期間有多筆季節性折扣時依序連續疊加。}{必要}
\end{reqtable}

\needspace{8\baselineskip}\textbf{報表統計子系統 (RSS)}
\begin{reqtable}
\req{RSS-F-001}{工作人員可以查詢每日銷售報表，內容包含日期、商店名稱、訂單數與銷售總額；已取消或已退貨的訂單不計入。}{必要}
\req{RSS-F-002}{工作人員可以查詢每週銷售報表（以週一至週日為一週）。}{必要}
\req{RSS-F-003}{工作人員可以查詢每月銷售報表。}{必要}
\req{RSS-F-004}{工作人員可以查詢各商店的暢銷商品報表。}{選配}
\req{RSS-F-005}{工作人員可以查詢指定期間內各商品或各商店的購買統計。}{選配}
\end{reqtable}

\needspace{8\baselineskip}\textbf{商店管理子系統 (STS)}
\begin{reqtable}
\req{STS-F-001}{訪客與顧客可以查看商店資訊（名稱、地址、聯絡 Email 與電話）與該商店販售的商品。}{必要}
\req{STS-F-002}{平台管理員可以修改商店資料。}{選配}
\req{STS-F-003}{顧客可以對已完成（closed）訂單中購買過商品的商店評分（1 至 5 分）並撰寫評論。}{選配}
\req{STS-F-004}{商店頁面顯示平均評分與評論列表。}{選配}
\end{reqtable}

\needspace{8\baselineskip}\textbf{資料庫子系統 (DBS)}
\begin{reqtable}
\req{DBS-F-001}{以關聯式資料庫（PostgreSQL）儲存系統所有資料，並實施 2.8 節的商業規則與完整性限制。}{必要}
\req{DBS-F-002}{資料庫初始狀態須包含平台管理員、工作人員與商店資料。}{必要}
\end{reqtable}

\subsection{資料需求 (Data Requirements)}
表 2 列出系統必須儲存的資料。詳細的實體、屬性、主鍵與關係將於 Phase 2 的 ER Diagram 與 schema 中定義。

\begin{xltabular}{\linewidth}{|>{\centering\arraybackslash}p{1.3cm}|>{\centering\arraybackslash}p{2.4cm}|X|}
\caption{資料需求}\\
\hline\rowcolor{headgray}\textbf{編號} & \textbf{資料項目} & \textbf{須儲存的資訊}\\\hline\endfirsthead
\hline\rowcolor{headgray}\textbf{編號} & \textbf{資料項目} & \textbf{須儲存的資訊}\\\hline\endhead
DR-01 & 使用者 & loginid、姓名、密碼（以雜湊值儲存）、Email、地址、電話、角色（顧客／工作人員／平台管理員，可多重）、帳號是否啟用、建立時間。\\\hline
DR-02 & 商店 & storeid、名稱、地址、聯絡 Email、電話。\\\hline
DR-03 & 商品 & productid、所屬商店、名稱、描述、分類、圖片、價格、庫存數量、是否上架、建立與更新時間。\\\hline
DR-04 & 購物車 & 顧客、商品、數量、加入時間。每位顧客的同一商品在購物車中只有一筆。\\\hline
DR-05 & 訂單 & orderid、下單時間、顧客、收件人姓名、電話與地址、商品小計、折扣總額、稅、運費、訂單總額、目前狀態（received、processing、shipping、closed、cancelled、returned）、付款確認時間。\\\hline
DR-06 & 訂單明細 & 所屬訂單、商品、下單當下的品名與單價、購買數量、明細小計、該項商品的特殊活動折抵金額。\\\hline
DR-07 & 訂單狀態歷程 & 訂單、狀態、變更時間、操作的工作人員。\\\hline
DR-08 & 折扣政策 & 折扣代碼、類型、描述、開始日期、結束日期，以及規則參數：折扣率、門檻金額、買 X 送 Y 的 X 與 Y、是否啟用。\\\hline
DR-09 & 特殊活動折扣適用商品 & 折扣代碼、商品，以及該折扣的期間（用於檢查同一商品期間不重疊）。\\\hline
DR-10 & 訂單套用折扣 & 訂單、折扣代碼、實際折抵金額。\\\hline
DR-11 & 系統參數 & 稅率、基本運費。\\\hline
DR-12 & 商店評價（選配） & 顧客、商店、訂單、評分、評論內容、評價時間。\\\hline
DR-13 & 退貨申請 & 申請編號、訂單、申請時間、退貨原因、處理狀態（pending、approved、rejected）、處理的工作人員、處理時間、拒絕原因、退款總額。\\\hline
DR-14 & 退貨明細 & 所屬退貨申請、訂單明細（商品）、退貨數量、該項退款金額。\\\hline
\end{xltabular}

銷售報表的內容（日期、商店名稱、訂單數、銷售總額）由上述訂單、訂單明細、商品與商店資料彙總產生，不另外儲存。訂單取消時不另外建立資料，只記錄在訂單狀態與狀態歷程中。

\subsection{非功能性需求 (Non-Functional Requirements)}

\subsubsection{效能需求 (Performance Requirements)}
\begin{reqtable}
\req{FES-N-101}{資料表的設計應避免重複儲存資料、節省空間。}{必要}
\req{FES-N-102}{在 1 萬筆商品的資料量下，關鍵字搜尋應於 2 秒內回傳結果。}{必要}
\req{FES-N-103}{一般頁面的讀取時間不應大於 3 秒。}{必要}
\req{FES-N-104}{一年份訂單資料下，月銷售報表應於 5 秒內產生。}{必要}
\req{FES-N-105}{多位顧客同時購買同一商品時，庫存不會被扣成負數（不超賣）。}{必要}
\end{reqtable}

\subsubsection{資安需求 (Security Requirements)}
\begin{reqtable}
\req{FES-N-201}{密碼以 bcrypt 雜湊後儲存，資料庫中不得出現明文密碼。}{必要}
\req{FES-N-202}{登入成功後以具有效期限的 JWT 識別使用者；每個 API 皆在伺服器端檢查角色權限，不只依靠前端隱藏按鈕。}{必要}
\req{FES-N-203}{所有 SQL 皆使用參數化查詢或 ORM，防止 SQL Injection。}{必要}
\req{FES-N-204}{訂單金額一律由伺服器計算，不接受前端傳入的價格與總額。}{必要}
\req{FES-N-205}{顧客只能讀取與修改自己的購物車、訂單與個人資料。}{必要}
\req{FES-N-206}{資料庫連線帳密等設定以環境變數管理，不放入版本控制。}{必要}
\end{reqtable}

\subsection{介面需求 (Interface Requirements)}

\subsubsection{使用者介面需求 (User Interfaces Requirements)}
\begin{reqtable}
\req{FES-N-301}{登入、登出與註冊介面}{必要}
\req{FES-N-302}{商品瀏覽與關鍵字搜尋介面}{必要}
\req{FES-N-303}{商品詳細資料介面}{必要}
\req{FES-N-304}{購物車介面}{必要}
\req{FES-N-305}{結帳與確認付款介面（顯示金額明細）}{必要}
\req{FES-N-306}{歷史訂單、訂單明細、取消訂單與退貨申請介面}{必要}
\req{FES-N-307}{工作人員商品管理介面}{必要}
\req{FES-N-308}{工作人員訂單處理與退貨審核介面}{必要}
\req{FES-N-309}{工作人員折扣政策管理介面}{必要}
\req{FES-N-310}{工作人員銷售報表介面}{必要}
\req{FES-N-311}{平台管理介面：使用者帳號、角色與權限管理，以及停用或啟用帳號}{必要}
\req{FES-N-312}{商店資訊與評價介面}{選配}
\req{FES-N-313}{導覽列依使用者角色顯示可用的功能；未授權的頁面無法進入。}{必要}
\req{FES-N-314}{介面支援繁體中文與英文切換，並支援桌機與手機瀏覽（RWD）。}{選配}
\end{reqtable}

\subsubsection{外部介面需求 (External Interface Requirements)}
\begin{reqtable}
\req{FES-N-401}{使用者以瀏覽器透過 HTTP(S) 與前端網頁伺服器通訊。}{必要}
\req{FES-N-402}{前端以 HTTP(S) 傳送 JSON 格式的 REST API 請求給後端 FastAPI 伺服器。}{必要}
\req{FES-N-403}{後端透過 SQLAlchemy 以 TCP 連線存取 PostgreSQL 資料庫。}{必要}
\req{FES-N-404}{後端透過 SMTP 寄送訂單確認 Email。}{選配}
\req{FES-N-405}{本系統不串接第三方金流；「確認付款」由顧客在結帳頁按下確認後即視為付款完成。}{必要}
\end{reqtable}

\subsubsection{內部介面需求 (Internal Interface Requirements)}
\begin{reqtable}
\req{FES-N-501}{AAS 和 DBS 間必須能傳送與接收使用者帳號、密碼雜湊與角色資訊。}{必要}
\req{FES-N-502}{AAS 和其他各子系統間必須能傳送目前登入使用者的身分與角色，供授權判斷。}{必要}
\req{FES-N-503}{PMS 和 DBS 間必須能傳送與接收商品資料與庫存。}{必要}
\req{FES-N-504}{SCS 和 DBS 間必須能傳送與接收顧客的購物車資訊。}{必要}
\req{FES-N-505}{SCS 和 OMS 間必須能傳送顧客欲結帳的商品與數量。}{必要}
\req{FES-N-506}{OMS 和 DMS 間必須能傳送訂單內容，並接收各項折扣的折抵金額。}{必要}
\req{FES-N-507}{OMS 和 PMS 間必須能在同一交易中檢查與扣除商品庫存，並在訂單取消或退貨核准時回補庫存。}{必要}
\req{FES-N-508}{OMS 和 DBS 間必須能傳送與接收訂單、訂單明細、狀態歷程、退貨申請與退貨明細。}{必要}
\req{FES-N-509}{DMS 和 DBS 間必須能傳送與接收折扣政策與適用商品資訊。}{必要}
\req{FES-N-510}{RSS 和 DBS 間必須能查詢訂單與商店資料以產生銷售報表。}{必要}
\req{FES-N-511}{STS 和 DBS 間必須能傳送與接收商店資料與評價資訊。}{必要}
\end{reqtable}

\subsection{其他需求 (Other Requirements)}

\subsubsection{環境需求 (Environmental Requirement)}
\begin{reqtable}
\req{FES-N-601}{開發與執行環境須同時支援 macOS（Apple Silicon）與 Windows 11。}{必要}
\req{FES-N-602}{後端使用 Python 3.12 以上版本，前端建置使用 Node.js 20 LTS 以上版本。}{必要}
\req{FES-N-603}{資料庫使用 PostgreSQL 16 以上版本。}{必要}
\req{FES-N-604}{使用者端支援最新版 Chrome、Edge 與 Safari，且需要網路連線。}{必要}
\end{reqtable}

\subsubsection{安裝需求 (Installation Requirement)}
\begin{reqtable}
\req{FES-N-701}{提供 Docker Compose 設定檔，可一鍵啟動 PostgreSQL 資料庫；亦可改用本機安裝的 PostgreSQL。}{必要}
\req{FES-N-702}{提供初始資料的匯入方式，安裝後即可得到含平台管理員、工作人員與商店資料的資料庫。}{必要}
\req{FES-N-703}{提供 README，說明安裝相依套件、設定環境變數與啟動前後端的步驟。}{必要}
\end{reqtable}

\subsubsection{測試需求 (Test Requirements)}
\begin{reqtable}
\req{FES-N-801}{每個子系統的功能都須經過測試。}{必要}
\req{FES-N-802}{驗證訂單金額計算正確，涵蓋三種折扣單獨套用、同時套用，以及多筆季節性折扣疊加的情況。}{必要}
\req{FES-N-803}{驗證 2.8 節的商業規則與完整性限制確實生效。}{必要}
\req{FES-N-804}{測試多位顧客同時購買同一商品最後幾件庫存時，不會發生超賣。}{必要}
\req{FES-N-805}{測試各角色只能使用被授權的功能。}{必要}
\req{FES-N-806}{於 FES-N-604 所列的瀏覽器上測試使用者介面。}{必要}
\req{FES-N-807}{測試取消訂單、全部與部分退貨的流程，以及庫存回補與退款金額計算是否正確。}{必要}
\end{reqtable}

\subsection{商業規則與限制 (Business Rules and Integrity Constrains)}

\begin{twocol}{編號}{規則描述}
\row{BR-01}{loginid、productid、storeid、orderid 與折扣代碼皆須唯一。}
\row{BR-02}{使用者分為顧客、工作人員與平台管理員三種角色，一位使用者可以同時具有多種角色（平台管理員可兼任工作人員）。初始資料庫至少包含一位平台管理員；工作人員與平台管理員角色只能由平台管理員指派或移除，且系統至少須保留一位啟用中的平台管理員。使用者帳號不提供刪除功能，只能停用（BR-20）。目前工作人員可以管理所有商店；各商店獨立管理權限屬於未來擴充。}
\row{BR-03}{每個商品屬於且僅屬於一間商店；一間商店可以販售零至多個商品。一筆訂單可以包含不同商店的商品。}
\row{BR-04}{商品價格須大於 0；庫存數量須為大於或等於 0 的整數。}
\row{BR-05}{每項購買商品的數量必須是大於 0 的整數（購物車與訂單明細皆適用）。}
\row{BR-06}{每項商品的購買數量必須小於或等於該商品目前的庫存數量；顧客確認付款後扣除庫存，訂單取消或退貨核准後回補庫存；庫存不得為負數。}
\row{BR-07}{每筆訂單屬於且僅屬於一位顧客，且至少包含一項商品；同一筆訂單中同一商品只出現一筆明細。}
\row{BR-08}{訂單明細保存下單當下的品名與單價；之後商品改名、改價或下架，不影響已成立的訂單。}
\row{BR-09}{訂單總額 = 商品小計 $-$ 特殊活動折扣 $-$ 季節性折扣 $+$ 稅 $+$ 運費 $-$ 運費折扣。金額以新台幣元為單位，稅額四捨五入至整數元。}
\row{BR-10}{金額計算順序：(1) 計算各商品小計；(2) 扣除特殊活動折扣（商品層級）；(3) 以扣除後金額計算季節性折扣（訂單層級）；(4) 以兩項折扣後的金額計算稅額，運費不課稅；(5) 若兩項折扣後的金額達運費折扣門檻則免運費。}
\row{BR-11}{運費折扣與季節性折扣作用於整筆訂單；特殊活動折扣作用於指定的商品。三種折扣可以同時套用於同一筆訂單。}
\row{BR-12}{同一期間有多筆季節性折扣有效時，全部依序連續疊加，例如 5\% 與 10\% 兩筆折扣後的金額為原金額 $\times 0.95 \times 0.90$。同一期間有多筆運費折扣時，只要達到任一筆的門檻即免運費。}
\row{BR-13}{同一商品在同一期間最多只能關聯一個特殊活動折扣；同一商品的各特殊活動折扣期間不得重疊（開始日與結束日皆包含在期間內）。}
\row{BR-14}{折扣期間的開始日期須小於或等於結束日期；折扣僅在期間內（含首尾兩日）有效，以顧客確認付款的日期判斷。}
\row{BR-15}{折扣率以小數儲存且須介於 0 與 1 之間（例如 0.05 代表打 95 折、折抵 5\%）。}
\row{BR-16}{「買 X 送 Y」表示每購買 X+Y 件同一商品只需支付 X 件的金額，X 與 Y 皆為正整數。}
\row{BR-17}{運費折扣的門檻金額須大於 0；季節性與特殊活動折扣須指定折扣率或買 X 送 Y 其中一種規則。}
\row{BR-18}{折扣代碼由系統產生並指派，建立後不可修改。}
\row{BR-19}{訂單狀態只能是 received、processing、shipping、closed、cancelled、returned 其中之一。正常流程只能依 received、processing、shipping、closed 的順序前進，不可跳過或倒退。只有 received 的訂單可以由顧客直接改為 cancelled；returned 只能在工作人員核准退貨申請、且訂單所有商品都已退貨時產生。部分商品退貨的訂單維持原狀態，繼續正常流程。cancelled 與 returned 為最終狀態。}
\row{BR-20}{已被訂單引用的商品、折扣政策與顧客帳號不可刪除，只能下架或停用，以維持參考完整性。}
\row{BR-21}{顧客只能存取自己的購物車、訂單與個人資料；只有工作人員可以管理商品、訂單狀態、退貨申請、折扣政策與查看報表；只有平台管理員可以管理使用者帳號、角色與權限。}
\row{BR-22}{稅率預設為 5\%，基本運費預設為新台幣 100 元，以系統參數儲存。}
\row{BR-23}{銷售報表中各商店的「銷售總額」以比例分攤計算：每筆訂單的訂單總額，依各商店商品扣除特殊活動折扣後的金額占全部商品的比例分攤給各商店，四捨五入至整數元後加總。「訂單數」為包含該商店商品的訂單筆數。狀態為 cancelled 或 returned 的訂單不計入報表；部分退貨的訂單，從退貨商品所屬商店的銷售總額扣除該項退款金額。}
\row{BR-24}{（選配）商店評分須為 1 至 5 的整數；顧客須有包含該商店商品且狀態為 closed 的訂單才能評價，每筆訂單對每間商店最多評價一次。}
\row{BR-25}{退貨申請只能針對狀態為 processing、shipping 或 closed 的訂單提出，可以只退部分商品；每項商品的退貨數量須為大於 0 的整數，且不得超過該商品的購買數量。每筆訂單只能提出一次退貨申請，被拒絕後不可再申請。訂單有待審核（pending）的退貨申請時，工作人員不可更新訂單狀態。}
\row{BR-26}{退款金額以比例分攤計算：每項退貨商品的退款金額 =（訂單總額 $-$ 運費）$\times$ 該項退貨商品扣除特殊活動折扣後的金額 $\div$ 全部商品扣除特殊活動折扣後的金額，四捨五入至整數元；退款總額為各項加總。運費不退，唯一例外是訂單在 processing（尚未出貨）時全部商品都退貨，運費一併退還。本系統不串接第三方金流，只記錄退款金額，實際退款不在系統範圍內。}
\end{twocol}

\clearpage
% ================= Glossary =================
\section*{Glossary}
\addcontentsline{toc}{section}{Glossary}

\begin{xltabular}{\linewidth}{|>{\centering\arraybackslash}p{3.4cm}|X|}
\hline\rowcolor{headgray}\textbf{名詞} & \textbf{說明}\\\hline\endhead
SRS & Software Requirements Specification，系統需求規格書。\\\hline
SPA & Single Page Application，單頁式網頁應用程式。\\\hline
REST API & 以 HTTP 方法（GET、POST、PUT、DELETE）操作資源的網路應用程式介面。\\\hline
JSON & JavaScript Object Notation，前後端交換資料使用的格式。\\\hline
JWT & JSON Web Token，登入後用來識別使用者身分的權杖。\\\hline
bcrypt & 一種加鹽的密碼雜湊演算法。\\\hline
ORM & Object-Relational Mapping，以物件操作關聯式資料庫的技術，本系統使用 SQLAlchemy。\\\hline
CRUD & Create、Read、Update、Delete，新增、查詢、修改、刪除。\\\hline
主鍵 (PK) & Primary Key，唯一識別資料表中每一筆資料的屬性。\\\hline
外鍵 (FK) & Foreign Key，參照其他資料表主鍵的屬性，用以維持參考完整性。\\\hline
交易 (Transaction) & 一組必須全部成功或全部失敗的資料庫操作。\\\hline
運費折扣 & 訂單金額達門檻時免運費的折扣，作用於整筆訂單。\\\hline
季節性折扣 & 指定期間內依折扣率折抵整筆訂單金額的折扣。\\\hline
特殊活動折扣 & 指定期間內作用於指定商品的折扣，規則為折扣率或買 X 送 Y。\\\hline
折扣代碼 & 系統指派給每個折扣政策的唯一代碼。\\\hline
received & 訂單已送出，系統已接收。\\\hline
processing & 工作人員處理中。\\\hline
shipping & 已出貨，配送中。\\\hline
closed & 訂單已完成。\\\hline
cancelled & 訂單在 received 狀態時由顧客取消。\\\hline
returned & 訂單所有商品的退貨申請都經工作人員核准，訂單已全部退貨。\\\hline
退貨申請 & 顧客對已處理、已出貨或已完成的訂單提出的退貨要求，可以只退部分商品，須經工作人員同意；每筆訂單限申請一次。\\\hline
RWD & Responsive Web Design，響應式網頁設計。\\\hline
\end{xltabular}

\clearpage
% ================= References =================
\section*{References}
\addcontentsline{toc}{section}{References}

\begin{enumerate}[label={[\arabic*]},leftmargin=2.5em,itemsep=0.4em]
  \item 國立臺北科技大學資訊工程系，資料庫系統課程專題說明（e-Shop Project Descriptions），2026。
  \item 國立臺北科技大學資訊工程系，系統需求規格書範本（SRS.docx）。
  \item NTUT Software Engineering Lab, Database System Example Project. \url{https://github.com/ntutselab/Database-System-Example-Project}
  \item IEEE Std 830-1998, IEEE Recommended Practice for Software Requirements Specifications.
  \item PostgreSQL Global Development Group, PostgreSQL 16 Documentation. \url{https://www.postgresql.org/docs/16/}
  \item FastAPI Documentation. \url{https://fastapi.tiangolo.com/}
  \item SQLAlchemy Documentation. \url{https://docs.sqlalchemy.org/}
  \item React Documentation. \url{https://react.dev/}
\end{enumerate}

\clearpage
% ================= Appendix =================
\section*{Appendix}
\addcontentsline{toc}{section}{Appendix}

無。

\end{document}
